\section{Conclusions and Future Work}
\label{sec:conclusion}

The results of the described system support three main conclusions.

First, analyzer and BM25 tuning provide a useful lexical baseline, especially thanks to its light computing print, with the best analyzer-focused run reaching 0.503 MRR@5 and 0.796 Recall@100 in the English run analysis.

Second, moving to BGE-M3 hybrid retrieval produces a stronger candidate generator: the retained configuration, \texttt{bge\_m3\_hybrid\_513}, reaches 0.557 MRR@5 and 0.839 Recall@100.

Third, reranking the hybrid candidates with \texttt{nvidia/llama-nemotron-rerank-1b-v2} gives the largest observed improvement among the reported single-stage runs, reaching 0.702 MRR@5 and 0.911 Recall@100 when reranking the top 1000 candidates. Alternative rerankers such as Qwen3-reranker-4B and \texttt{Alibaba-NLP/gte-reranker-modernbert-base} also improve over retrieval alone, but remain below the Nemotron runs in the reported MRR@5 values.

The development results show that fine-tuning the Nemotron reranker is promising. Fine-tuning improves the performance by a small but appreciable amount (average of +0.03 of MRR@5 in our experiments). Nevertheless we recognize the limits of our knowledge, so a higher-quality fine-tuning, not only for the reranker, can bring better results is certainly possible.

\subsection{Limitations}

Our system exhibits several important limitations. First, performance degrades significantly for non-English queries: German achieves only 0.674 MRR@5 compared to French's 0.760 and English's 0.716 on the development set. This gap is attributable to three factors: (1) the German development set is small (only 386 labeled queries), limiting statistical confidence; (2) machine translation from German to English may introduce morphological errors that confuse the retrieval and reranking models; and (3) German's complex morphology and abundant compound words may require specialized handling not captured by the standard translation pipeline.

Second, our system relies exclusively on title and abstract representations, which may omit important evidence from authors, publication venue, or full-text content. This design choice was made for efficiency but may miss relevant signals for resolving ambiguous claims.

Third, the hard-negative mining strategy is limited to five negatives per query within a similarity range of $[0.05, 0.85]$ cosine distance from the positive. More sophisticated strategies such as curriculum learning or loss-weighted hard-negative sampling might yield further improvements.

Fourth, our cascaded reranking experiments (Table~\ref{tab:aggregate-rerankers}) show promise but lack complete per-query outputs for rigorous comparative analysis. These results are therefore presented as aggregate diagnostics only and would benefit from full reproducibility.

\subsection{Future Work}

Several directions emerge for future improvement. First, deploying a fine-tuned German translation model could mitigate morphological drift and improve the performance of non-English queries. Second, implementing and validating the cascaded reranking configuration documented in Table~\ref{tab:aggregate-rerankers} could yield the highest observed MRR@5 (0.766) on English data, though at increased computational cost.

Third, in-depth analysis of the 929 completely missing queries (6.2\% of the English split) could reveal systematic failure modes—such as claims referencing methodology or findings not mentioned in paper abstracts, or highly paraphrased social-media versions that bear little lexical resemblance to the original publication.

Fourth, incorporating additional document representations such as author lists, publication venues, keywords, or full-text content could enable richer retrieval and reranking signals for claims mentioning specific research groups, conferences, or specialized domains.

Fifth, ensemble approaches that combine predictions from multiple independently fine-tuned rerankers with learned fusion weights might improve robustness and generalization across the multilingual test set.

Finally, analyzing the relationship between query characteristics (length, language, entity density, paraphrase ratio) and retrieval difficulty could inform targeted improvements for challenging query subsets.

\section {Declaration on Generative AI}
During the writing of this paper the authors used Grammarly and DeepL for grammar and spelling checks and Claude Sonnet 4.6 for writing-style improvement, citation management and formatting assistance. After using these tools, the authors reviewed and edited the content as needed and take full responsibility for the publications content.